\section{The cubic Type II block: an integral descent and its remaining vanishing input}
\label{sec:cubic-type-ii-descent}

The rank-three deformation argument has a different coefficient interface
from the one-block theta argument.  In this section the quadratic field
is $K=\mathbb Q(\omega)$, $\omega^3=1$, $\omega\ne1$, and
$\chi=\chi_{269}$ is the cubic character of $K(\sqrt[3]{269})/K$.
The following descent repairs one precise part of a possible cubic
version of the rank-three argument.  It does not assert that its new
deformation problem satisfies the supplied automorphy or rigidity inputs.

\begin{proposition}[AC-G45: the actual group and integral Type II descent]
\label{AC-G45}
Let $\Gamma=H\rtimes\langle c\rangle$, with $c^2=1$, and let
$\epsilon:\Gamma\longrightarrow\{\pm1\}$ have kernel $H$.
Let $R$ be a commutative ring, let
\[
 J_3=\begin{pmatrix}0&0&1\\0&1&0\\1&0&0\end{pmatrix},
 \qquad \vartheta(A)=J_3 A^{-t}J_3,
 \qquad \mathcal G_3(R)=\operatorname{GL}_3(R)
                      \rtimes_{\vartheta}\{\pm1\}.       \tag{G45.1}
\]
Suppose $r:\Gamma\to\mathcal G_3(R)$ has second component
$\epsilon$, has $r(c)=(I,-1)$, and is block diagonal on $H$
for the decomposition $R\{e_0,e_2\}\oplus Re_1$.
Write its endpoint block as $B(h)\in\operatorname{GL}_2(R)$.
Let $\xi:H\to R^\times$ satisfy $\xi(chc)=\xi(h)^{-1}$,
and suppose $\det B(h)=\xi(h)^2$.
Then the representation $\xi^{-1}B$ of $H$ has an explicit
integral extension $V:\Gamma\to\operatorname{GL}_2(R)$:
\[
 C(h)=\xi(h)^{-1}B(h),\qquad
 D=\operatorname{diag}(-1,1),\qquad
 V(hc^a)=C(h)D^a\quad(a=0,1).                         \tag{G45.2}
\]
Its determinant is $\epsilon$, and $\operatorname{tr}V(c)=0$.
Continuity and irreducibility on $H$ are preserved whenever
these notions are applicable.  No division by $2$ or $3$, no
Fourier projector, and no extraction of a square root is required.

In particular, a cubic Type II point whose middle character is $1$
and whose endpoint determinant is $\chi^2$ has the form
\[
                 \rho_K\simeq\chi\,V|_{G_K}\oplus1,   \tag{G45.3}
\]
where $V$ is odd and has determinant $\eta_K$.
The middle character need not equal $1$ for the descent (G45.2)
itself: after the scalar untwist it is allowed to be $\chi^{-1}$.
\end{proposition}
\begin{proof}
\proofstep{1. Exact supplied definitions.}
All source locators here use commit
\texttt{7cbf7a0c32fca1badf05931879a37c3f8a2b03b4} and begin with
\path{Development/AutomorphicArithmetic/AutomorphicVarieties/}.
In \path{UnitaryThreeGroups/UnitaryThreeReflectionRight.lean},
\texttt{Jmat}, \texttt{theta}, \texttt{sgnAct}, and \texttt{GG}
define exactly (G45.1).  Thus the ambient group is a general-linear
semidirect product with inverse-transpose action; its definition
does not impose determinant one on the three-dimensional matrix.
The same file defines
\texttt{r0} by $\operatorname{diag}(\Psi,1,\Psi^{-1})$
and defines \texttt{Normalised.atc} by $r(c).\mathrm{mat}=I$.

In \path{UnitaryThreeGroups/BaseRepGL3NormalisedOfPsiEqOne.lean},
\texttt{Block}, \texttt{corner}, \texttt{TypeII}, and \texttt{Wfun}
are respectively the vanishing of the odd-weight entries, the
submatrix on indices $\{0,2\}$, irreducibility of that submatrix
on $\ker\eta$, and multiplication of the corner by $D$ on the
nontrivial quadratic coset.  Type II therefore uses the two
endpoint coordinates, not a middle--endpoint pair.
The source theorem
\path{UnitaryThreeGroups/ExistsLieElementMiddleBlockPart.lean},
\texttt{exists\_lieElement\_middleBlockPart}, constructs this
extension under the endpoint determinant-one hypothesis.
The full proof, including \texttt{corner\_theta'}, uses no assertion
that the middle entry is $1$.
The same distinction is explicit in
\path{UnitaryThreeGroups/CoboundaryLemmaNoColumn.lean},
private \texttt{W\_facts}, lines 101--133, and
\texttt{typeII\_facts}, lines 135--215: the latter first proves
the middle and determinant identities and then invokes the former.

\proofstep{2. Scalar untwisting is defined on the semidirect product.}
Put $a(hc^u)=\xi(h)$.  The hypothesis on $\xi$ gives
\[
             a(gg')=a(g)a(g')^{\epsilon(g)}.          \tag{G45.4}
\]
Since $\vartheta(zI)=z^{-1}I$, replacing the matrix component
$A(g)$ of $r(g)$ by $a(g)^{-1}A(g)$ defines another homomorphism
to $\mathcal G_3(R)$.  This follows by substituting (G45.4)
in the semidirect product multiplication law.  The replacement
fixes the matrix at $c$, preserves the endpoint/middle block
decomposition, and makes the endpoint determinant equal to $1$.
It is a scalar twist of the whole three-dimensional matrix, so
its effect on the middle entry is retained rather than discarded.

\proofstep{3. The endpoint extension.}
Let $J_2=\left(\begin{smallmatrix}0&1\\1&0\end{smallmatrix}\right)$.
The relation $r(chc)=r(c)r(h)r(c)$ gives, after untwisting,
\[
 C(chc)=J_2C(h)^{-t}J_2.
\]
For $C=\left(\begin{smallmatrix}a&b\\d&e\end{smallmatrix}\right)$
with $ae-bd=1$, direct inversion gives
\[
       J_2C^{-t}J_2=
       \begin{pmatrix}a&-b\\-d&e\end{pmatrix}=DCD.
\]
Consequently $C(chc)=DC(h)D$ and $D^2=I$, which are exactly
the relations needed for (G45.2) to be a homomorphism.
Its determinant is $1$ on $H$ and $-1$ on the other coset;
its trace at $c$ is $-1+1=0$.
The formulas involve only the given integral matrices and unit
character values.  A subspace invariant under $V$ is invariant
under $C(H)$, so irreducibility on $H$ implies irreducibility
of $V$.  Continuity follows on each of the two open cosets.
For $\Gamma=G_{\mathbb Q}$ and a complex conjugation $c$,
$H=G_K$ and $\epsilon=\eta_K$, proving (G45.3).

The symmetric inverse-transpose form is essential to this oddness
conclusion.  If the conjugate-dual form on a two-dimensional
determinant-one representation were instead the alternating matrix
$J=\left(\begin{smallmatrix}0&1\\-1&0\end{smallmatrix}\right)$,
then $J C^{-t}J^{-1}=C$, so this argument would give a commuting
involution, not $D$.  For an irreducible representation over an
algebraically closed field, such a commuting involution is scalar
and therefore even.  The actual source form is the symmetric
anti-diagonal $J_2$, so this ambiguity does not occur here.
\end{proof}

\begin{proposition}[AC-G45-local: a sufficient cubic local condition]
\label{AC-G45-local}
Let $\ell\notin\{2,3,269\}$, and let $\mathcal O$ be a finite
extension DVR of $\mathbb Z_\ell$ containing the character values.
Let $\Psi=\psi\kappa^e$ be a character of $G_{\mathbb Q}$,
where $\psi$ is finite and unramified at $269$, and $\kappa$
is cyclotomic.  Write $\gamma=\Psi|_{G_K}$ and put
\[
        \rho_0(h)=\operatorname{diag}
                    (\chi(h)\gamma(h),1,\chi(h)\gamma(h)^{-1}).
                                                               \tag{G45.5}
\]
Suppose an integral continuous block representation
$\rho:G_K\to\operatorname{GL}_3(\mathcal O)$ satisfies:
\begin{enumerate}
\item $\overline\rho=\overline\rho_0$ in the specified basis;
\item for every finite inertia group $I_v$ and $h\in I_v$,
      $\det(X-\rho(h))=\det(X-\rho_0(h))$;
\item for $v\nmid\ell$ and $h\in I_v$ with
      $\overline\rho_0(h)=I$, one has $\rho(h)=I$;
\item at each $v\mid\ell$ there is $s\in I_v$ with
      $\overline\gamma(s)\ne1$.
\end{enumerate}
Then the middle character of $\rho$ is $1$, and its endpoint
determinant is $\chi^2$.  If $\rho$ comes from a normalized
homomorphism to (G45.1), AC-G45 applies integrally.

These are sufficient hypotheses for a modified deformation problem,
not assertions that the existing \texttt{Dat} has already been
modified or that its automorphy theorem accepts this new seed.
\end{proposition}
\begin{proof}
The field and ramification calculations in AC-K04 give that
$\chi$ is unramified away from $269$ and is trivial at the
prime above $3$.  In particular it is unramified at $\ell$.
The ring $\mathbb Z[\omega]$ is norm-Euclidean and $K$ has
class number one, as used in AC-K12.
Let $m$ be the middle character of $\rho$.  Its reduction is $1$.
For $h\in I_v$, hypothesis 2 implies that $m(h)$ is one of
the three displayed eigenvalues of $\rho_0(h)$.

If $v\nmid269\ell$, then $\chi(h)=1$.  When
$\overline\gamma(h)\ne1$, both $\gamma(h)$ and its inverse
have reduction different from $1$, so $m(h)=1$.
When $\overline\gamma(h)=1$, hypothesis 3 gives the same
conclusion.  If $v\mid269$, then $\gamma(h)=1$ because
$\psi$ and $\kappa$ are unramified there.
The nontrivial cubic roots remain different from $1$ in residue
characteristic $\ell\ne3$: indeed $1+\omega+\omega^2=0$
would otherwise reduce to $3=0$.
Thus $\chi(h)\ne1$ again singles out $m(h)=1$; if
$\chi(h)=1$, hypothesis 3 applies.

Finally let $v\mid\ell$.  Again $\chi|_{I_v}=1$.
The same eigenvalue argument works when
$\overline\gamma(h)\ne1$.  Otherwise choose $s$ from
hypothesis 4.  Both $s$ and $hs$ have nontrivial residual
$\gamma$-value, so $m(h)=m(hs)m(s)^{-1}=1$.
This proves that $m$ is unramified at every finite prime.
Every continuous unramified abelian character of the imaginary
quadratic field $K$ factors through its ideal class group.
Since that group is trivial, $m=1$ globally.
This is exactly the class-field-theoretic mechanism of the supplied
\path{UnitaryThreeGroups/PowEqOneOfUnr.lean},
\texttt{pow\_eq\_one\_of\_unr} and \texttt{char\_eq\_one};
here the class number is $1$, so no further near-identity argument
is needed after unramifiedness has been proved.

The determinant identity in hypothesis 2 says
$\det\rho(h)=\chi(h)^2$ on every finite inertia group.
Thus $\det\rho\,\chi^{-2}$ is another everywhere unramified
character and is trivial for the same reason.  Since the middle
character is $1$, the endpoint determinant is $\chi^2$.
\end{proof}

\begin{proposition}[AC-G45-cohomology: the new rigidity coefficient]
\label{AC-G45-cohomology}
Under (G45.3), the part of the semidirect adjoint representation
consisting of the four off-diagonal endpoint/middle entries is
\[
 \mathcal R\simeq
 \operatorname{Ind}_{G_K}^{G_{\mathbb Q}}(\chi\,V|_{G_K})
 \simeq V\otimes\operatorname{Ind}_{G_K}^{G_{\mathbb Q}}\chi.
                                                               \tag{G45.6}
\]
The diagonal-block adjoint action is unchanged by the scalar
character on the endpoint block.  When $\chi=1$, (G45.6) is
$V\oplus V\eta_K$; when $\chi$ has order three,
$\operatorname{Ind}\chi$ is irreducible of dimension two and
has determinant $\eta_K$.

For coefficient-field representations, induction identifies the
Bloch--Kato Selmer group of $\chi V|_{G_K}$ over $K$ with that
of $\mathcal R$ over $\mathbb Q$.  Consequently the new Type II
off-diagonal vanishing input is the rank-four assertion
\[
 H^1_f\!\left(\mathbb Q,
        V\otimes\operatorname{Ind}_{G_K}^{G_{\mathbb Q}}\chi
        \right)=0,                                      \tag{G45.7}
\]
with the source's Tate normalization of $V$.
The presently supplied \texttt{KatoUse} and \texttt{KatoTwoW}
do not state (G45.7).  Odd endpoint descent alone does not
establish modularity of all the new points, (G45.7), or a
quantitative cubic Selmer bound.
\end{proposition}
\begin{proof}
Write $B=\chi V|_{G_K}$.  On $G_K$ the column from the
middle line to the endpoint block transforms as $B$, while
the opposite row transforms as $B^\vee$.  The source adjoint
formula in \texttt{UnitaryThreeReflectionRight.lean},
\texttt{thetaLie}, \texttt{GG.adFun}, and \texttt{adRep},
makes $c$ exchange these two two-dimensional spaces by negative
transpose and the anti-diagonal matrices.  Its square is $1$.
The first column space and its $c$-translate therefore give the
two coset summands in the induced representation.  Explicitly,
send the identity-coset copy of $B$ to that column space and
the other copy to its image under $c$.  The action of $G_K$
and of $c$ agrees on both copies, proving the first isomorphism
in (G45.6).  Tensor induction gives the second: on each of the
two cosets use the already defined extension $V$ to transport
the vector factor.  The endpoint endomorphism action is
$(\chi V)X(\chi V)^{-1}=VXV^{-1}$; the middle scalar and
the action of $c$ are also unchanged.  This proves the assertion
about the other blocks.

For $\chi=1$, the two coset permutation representation has
the invariant and anti-invariant lines, giving $1\oplus\eta_K$
over the coefficient field.  For nontrivial cubic $\chi$, its
restriction to $G_K$ is $\chi\oplus\chi^{-1}$ with distinct
characters.  A one-dimensional invariant subspace would have
to be one of these eigenspaces, but $c$ exchanges them.
The induced representation is therefore irreducible.  Its
determinant is $1$ on $G_K$ and $-1$ at $c$, as claimed.

For completeness, the Selmer comparison here is a comparison
of the specified local conditions as well as of global $H^1$.
The finite-coset cochain isomorphism for induction commutes
with restriction to a decomposition group; there it is the
direct sum over the primes of $K$ above the given prime of
$\mathbb Q$.  At a prime away from $\ell$ it carries the
kernel of restriction to inertia to the corresponding direct
sum of kernels.  At $\ell$, use the same finite-coset map
after tensoring with $B_{\mathrm{cris}}$: the elementary
tensor-induction identity commutes with the coefficient map
to $B_{\mathrm{cris}}$, so the kernels defining $H^1_f$ also
correspond.  Taking the kernel of the product of local maps
proves the assertion.  This statement concerns coefficient
fields; it is not an assertion that an integral projector at
the prime $2$ is harmless.

Finally, the exact source domain is explicit.
\path{UnitaryThreeGroups/KatoUse.lean}, \texttt{KatoUse},
lines 13--17, quantifies over $W:G_{\mathbb Q}\to\mathrm{GL}_2$
attached, after its Tate normalization, to a newform of weight
$2e+1$, and over a finite character of $G_{\mathbb Q}$.
It concludes vanishing for a two-dimensional character twist.
The newer
\path{ModularityArgument/KatoQuadraticCharacters/SphericalDecomposition/KatoTwoW.lean},
\texttt{KatoTwoW}, lines 11--15, still quantifies over
$W:G_{\mathbb Q}\to\mathrm{GL}_2$ and characters
$\theta,\psi$ of $G_{\mathbb Q}$, and concludes the two
vanishings for $W$ and $W\theta$.
The proof of
\path{ModularityArgument/ThresholdArguments/DenseThresholdCases/FPtHypU.lean},
\texttt{ptHyp\_of\_autB\_of\_lemW\_of\_isBase},
uses exactly that Type II vanishing to deduce rigidity of the
full adjoint representation.
More explicitly, the same directory's
\path{TypeIIOfBlochKatoLocalPoint.lean},
\texttt{bloch\_kato\_selmer\_vanishes\_of\_type\_ii\_modular\_ext},
sets $\theta=\eta_K$, invokes \texttt{kUse'} for $W$ and
$W\theta$, and uses the maps \texttt{pB 1},
\texttt{pB (-1)}, \texttt{sB 1}, and \texttt{sB (-1)}
to split and recombine the four off-diagonal entries.
Its final recombination uses the coefficient $1/2$.
Thus (G45.6) identifies the exact source block whose two-character
decomposition is changed by a cubic twist.
The irreducible two-dimensional factor $\operatorname{Ind}\chi$
is not one of those one-dimensional characters.
Furthermore \texttt{Dat.IsExt} and \texttt{Dat.IsModular}
in \path{UnitaryThreeGroups/RhoKApply.lean}, lines 85--92,
and \texttt{LemW} in
\path{AutomorphicCases/AutomorphyUnramifiedCoefficients.lean},
lines 18--22, are quantified over the original rational-character
datum.  The explicit descent just proved supplies the required
kind of odd two-dimensional representation but does not transport
those universally quantified automorphy hypotheses to a new datum.
\end{proof}

\begin{remark}[Scope of the remaining rank-three route]
The new content is the source-normalized integral descent and
the identification of the precise additional cohomology block.
At a rational prime inert in $K$, the unramified cubic character
has trivial value at the prime of $K$; this is compatible with
the unchanged Frobenius comparator studied separately in AC-K29.
The comparator does not construct the new deformation points.
A cubic version must still prove its local deformation and
global extraction statements, modularity of the Type I and
Type II points with the required level bounds, and the new
Type II vanishing (G45.7).  The exceptional coefficient primes
excluded in AC-G45-local are not covered by that proposition.
No claim is made that a generic Rankin--Selberg Euler-system
theorem automatically proves (G45.7).
\end{remark}

\paragraph{Geometric dictionary.}
After removing the finite ramification set, the finite cubic
character is the rank-one lisse coefficient associated with
its finite \'etale cover.  The endpoint tensor operation and
the induction in (G45.6) are the corresponding operations on
representations of the \'etale fundamental group.  This uses
the already audited dictionary of SGA~1, Expos\'e~V,
\S\S4 and~7.  The cubic subcover is the Kummer torsor of
SGA~1, Expos\'e~XI, \S6, Propositions~1 and~4 and
Corollary~5, after inverting $3\cdot269$.
These are categorical interpretations of the explicit descent;
they assert neither a new archival reading nor a geometric
vanishing theorem for the rank-four Bloch--Kato coefficient.
